\section{Design of \sys}
\label{sec:design}

In this section, we describe the design of \sys, which assesses vulnerability exploitability through downstream applications by recovering the calls and runtime state needed to trigger vulnerable behavior and finding inputs that trigger it.

\PP{Overview.}
\autoref{fig:overview} shows the workflow of \sys, which consists of three phases: \textit{trigger path recovery} (P1), \textit{harness generation} (P2), and \textit{concolic execution} (P3).

\begin{figure}[t]
\centering
\includegraphics[width=\linewidth]{./design/overview_HJ_simple.pdf}
\caption{Overview of \sys.}
\label{fig:overview}
\end{figure}

In P1, \sys recovers a call sequence from a public API of the downstream application to the vulnerable code in the upstream library, along with the conditions required to reach it.
In P2, \sys generates a harness that invokes the downstream application's public APIs in the recovered sequence and establishes the required arguments and runtime state.
Together, P1 and P2 recover how the downstream application can reach the vulnerable code and prepare the state needed to trigger it.
In P3, \sys performs concolic execution on the generated harness to find inputs that trigger the vulnerability.


\subsection{Recovering Trigger Path (P1)}
\label{subsec:p1}

The goal of P1 is to recover how a client of the downstream application $D$ can reach a sink $s$ in the upstream library $U$.
A client never calls $U$ by name.
It may invoke only $D$'s interface $I$---the functions that $D$ both exports and declares in its headers---and control reaches $U$ only after $D$ has installed the callees that cross the boundary.
What P1 recovers is therefore a \emph{trigger path}: a call sequence
\[
\pi = f_0 \rightarrow f_1 \rightarrow \cdots \rightarrow f_n
\]
together with a path condition $\varphi$ under which that sequence can run.
Here $f_0 \in I$, each step $f_i \rightarrow f_{i+1}$ is a resolved call, and $f_n$ contains $s$.
P2 establishes $\varphi$ through $D$'s APIs, and P3 searches for concrete bytes on the harness that P2 builds.
P1 itself is a \emph{may} analysis: a retained $(\pi, \varphi)$ is a candidate that install-resolved edges and a satisfiable condition at $f_0$ make plausible, not a proof that some client execution reaches $s$.
Spurious candidates are discarded when P2 cannot establish $\varphi$ or when P3 finds no concrete witness.

Tracking $\pi$ takes three analyses.
The call graph must expose the steps $f_i \rightarrow f_{i+1}$, including indirect ones.
The start $f_0$ must be a client-visible root in $I$.
The sequence must then survive a check that $\varphi$ is satisfiable at $f_0$.
Indirect calls dominate the $D$--$U$ boundary, so a graph of direct edges alone misses real steps, while signature matching invents steps that no install wrote.
Even on a careful graph, $\pi$ may still be unrealizable because of branches or missing installs.
P1 answers the first issue with field-sensitive points-to analysis of function-pointer installs, and the second by solving $\varphi$ at $f_0$.

The input is closed LLVM bitcode for $D$ linked with $U$, plus the source location of $s$, mapped to an IR instruction by debug information.
$I$ is the intersection of dynamically exported symbols and header-declared names.

\PP{Phase overview.}
P1 recovers $\pi$ in order.
Call-target resolution builds the edges $f \rightarrow g$ of the call graph~(\autoref{subsubsec:resolve}).
Entry extraction chooses admissible starts $f_0$~(\autoref{subsubsec:entries}).
Path enumeration extends $f_0$ toward $s$ into a full sequence and retains $(\pi, \varphi)$ when $\varphi$ is satisfiable at $f_0$~(\autoref{subsubsec:enumerate}).
The resolved edges are reused by P2 and P3, so later phases do not re-analyze calls.

\subsubsection{Resolving Call Targets}
\label{subsubsec:resolve}

The first analysis recovers the steps of $\pi$: for every call site in a function $f$, which callees $g$ that site may enter, so that each consecutive pair in $\pi$ is an edge $f \rightarrow g$ on the call graph.
Direct calls name $g$ in the IR and contribute that edge immediately.
Indirect calls do not: they load a location and jump to the function constant stored there.
Downstream code often reaches upstream code through exactly those sites---a setup path stores a handler into an object field, and a later path loads that field---so the possible $g$'s are what this analysis must compute.
Ignoring the site drops a real step of $\pi$.
Filling it with every address-taken function of matching type invents steps that no store created.
P1 therefore computes $g$ from pointer installs.

\PP{Points-to analysis of function-pointer installs.}
P1 runs a field-sensitive, flow-insensitive points-to analysis restricted to function pointers~\cite{lu2019where,cheng2000access}.
It tracks each function constant $g$ defined in $D \cup U$ through the abstract locations that may hold it.
A location $\ell$ is a root (a global, an argument, or a heap object keyed by allocation site) refined by an access path of struct-field indices and constant offsets when the layout is known---so a store into one field never updates another---and the same paths are recovered from TBAA tags when IR uses byte GEPs.
Flow insensitivity means one tracked set per $\ell$ for the whole module: every IR transfer that moves a function pointer updates that set, and the analysis keeps the least fixed point of those updates.

Tracking begins at stores of a function constant: if $g$ is written into $\ell$, $g$ is tracked at $\ell$.
A cast, phi, select, or GEP from $u$ to $v$ carries every tracked $g$ at $u$ over to $v$.
A store of a pointer into address $a$ carries the tracked set into every cell of $a$, and a load carries those cells' sets into the loaded value; tokens also move between a field path and every suffix of that path, so the same field is shared under outer and inner addressing.
At a call, tracked values move from actuals into formals of every type-compatible callee already tracked at the callee operand, and from that callee's return into the call result; as the operand's tracked set grows, new callees are admitted until the fixed point.

At a non-virtual indirect call in $f$, P1 reads the tracked set of the loaded location and keeps every $g$ from $D \cup U$ whose type matches the call.
Each such $g$ becomes an edge $f \rightarrow g$.
If nothing is tracked there, the site contributes no library edge: P1 does not invent steps from signature matches alone.

Virtual calls are not resolved by the same tracking.
Untyped object contents erase field offsets, so a vptr cannot share that cell with ordinary data without merging unrelated values.
P1 therefore takes vtable methods of classes derived from the receiver's static type (from debug inheritance) and retains those that match the call type.

\PP{Completing the call-graph edges.}
The tracked installs above recover stores performed inside $D$ and $U$.
The same loaded location may still receive a function pointer from outside the closed bitcode---a client callback or a runtime hook---so P1 also records those external callees and treats an edge that exists only that way as an obligation for the harness or the environment.
The edges used when extending $\pi$ are the tracked callees together with those external ones, restricted to address-taken functions of matching type.
P2 and P3 read this same edge relation and do not re-resolve calls.

\subsubsection{Selecting Entry Points}
\label{subsubsec:entries}

The call graph supplies possible steps $f_i \rightarrow f_{i+1}$, but not yet the first function $f_0$ of $\pi$.
P1 next chooses that start.

It performs a reverse reachability search from the function that contains $s$, walking edges backward, and obtains the cone of every function that can reach $s$.
Intersecting the cone with $I$ yields the exports that can appear on some path to the sink.
Some of those exports appear in the cone only because another interface function already called them.
They can be an $f_i$ for $i>0$, but not $f_0$: treating them as starts would force the harness to invent a prefix that $\pi$ never records.
P1 therefore keeps only cone roots inside $I$---functions $f_0 \in I$ that reach $s$ and have no caller inside the cone.
Every retained $\pi$ begins at such an $f_0$.

\subsubsection{Enumerating Feasible Paths}
\label{subsubsec:enumerate}

P1 finally builds full sequences $\pi = f_0 \rightarrow \cdots \rightarrow f_n$ on the call graph, where $f_n$ contains $s$.
From each admissible $f_0$ it searches for simple paths to $f_n$, ranking candidates by fewer indirect edges and then by shorter length.
After a path is accepted, its last edge into $f_n$ is banned and the search resumes, so later paths enter the sink at a different call site.

A path returned by this search is still only a syntactic $\pi$.
Inside some $f_i$, branches may refuse the route to the call of $f_{i+1}$, and an indirect step $f_i \rightarrow f_{i+1}$ succeeds only when the loaded location already holds $f_{i+1}$.
P1 therefore materializes a path condition $\varphi$ for each candidate and retains $(\pi, \varphi)$ only when a solver finds a model at $f_0$.

\PP{Forming and checking the path condition.}
For a candidate $\pi$, P1 builds the intraprocedural CFG of each $f_i$, collects the basic blocks that can reach the call to $f_{i+1}$ (or the sink instruction when $i=n$), and selects one short route from $f_i$'s entry across that slice.
Branch predicates on those routes become control constraints.
Each indirect step $f_i \rightarrow f_{i+1}$ contributes a dispatch constraint that the loaded location holds $f_{i+1}$.
$\varphi$ is the conjunction of those control and dispatch constraints.
P1 asks an SMT solver for a model of $\varphi$ over the arguments and fields of $f_0$, following stores into locals and summarizing callees that only write constant offsets into export-visible objects.
If $\varphi$ is satisfiable at $f_0$, P1 retains $(\pi, \varphi)$ for P2 and P3.
Otherwise the candidate is discarded, while another route through the same functions may still succeed.

\PP{Output of P1.}
P1 emits the retained pairs $(\pi, \varphi)$ and the resolved call-graph edges.
If nothing is retained, no sequence $f_0 \rightarrow \cdots \rightarrow f_n$ with $f_0 \in I$ explains how a client of $D$ reaches $s$.
A non-empty result is still only a may set of trigger paths for P2 and P3 to discharge.


\subsection{Generating Harness (P2)}
\label{subsec:p2}


\subsection{Concolic Execution (P3)}
\label{subsec:p3}
